\section{The arithmetic trace quotient and exact heat transport}
\label{sec:connes-complete-transport}

\subsection{The primary arithmetic object and the topology actually specified}
\label{sec:cq-source}
The starting point is Connes--Consani--Moscovici,
\emph{Zeta zeros and prolate wave operators: Semilocal adelic operators},
Section 3.6, proof-PDF pages 17--19, especially Lemma 3.3,
equations (3.16)--(3.22), and Proposition 3.6
\cite{ccmProlate2024}. The preceding arithmetic map is
\[
 \mathcal E f(x)=x^{1/2}\sum_{n\geq1}f(nx),\qquad
 \mathcal F_\mu g(s)=\int_0^\infty g(x)x^{-is}\,\frac{dx}{x}.
\]
The input consists of even Schwartz functions with
$f(0)=\widehat f(0)=0$. The paper's two Hermite families are
\[
 \psi_\ell^+=h_{4\ell}-\frac{h_{4\ell}(0)}{h_0(0)}h_0,
 \qquad
 \psi_\ell^-=-h_{4\ell+2}
       +\frac{h_{4\ell+2}(0)}{h_2(0)}h_2\quad(\ell\geq1).
\]
These signs and the minus sign in the Mellin exponent are retained.
The completed function is
\begin{equation}\label{eq:cq-Xi}
 \Xi(s)=\xi(\tfrac12+is),\qquad
 \xi(z)=\tfrac12z(z-1)\pi^{-z/2}\Gamma(z/2)\zeta(z),
 \qquad \Xi(s)=8H_0(2s).
\end{equation}
The paper proves
\begin{equation}\label{eq:cq-image}
 \mathcal F_\mu\mathcal E\psi_\ell^\pm=R_\ell^\pm\Xi,
 \qquad P_\ell^\pm=-\tfrac12(\tfrac14+s^2)R_\ell^\pm.
\end{equation}
It also states immediately before Proposition 3.6 that these give all
polynomial multiples of $\Xi$. Here is the precise algebra behind the
last assertion. The top summands of (3.19) and (3.20) are nonzero;
their degrees are $2\ell$ and $2\ell+1$. Division by
$-\tfrac12(\tfrac14+s^2)$ gives degrees $2\ell-2$ and
$2\ell-1$, respectively. The even and odd parities remain.
Thus the ordered family $R_1^+,R_1^-,R_2^+,R_2^-,\ldots$
has one polynomial of each degree, with a nonzero leading coefficient.
Subtracting the appropriate multiple of the polynomial of largest
degree, and iterating down to degree zero, expresses every polynomial
as a finite linear combination. Linear independence follows by the
largest-degree coefficient of a finite relation. Consequently
\begin{equation}\label{eq:cq-N}
 \cqN=\overline{\operatorname{span}
 \{\mathcal F_\mu\mathcal E\psi_\ell^\pm\}}^{\,\tau}
 =\overline{\cqC[s]\Xi}^{\,\tau},\qquad
 \cqQ=(\cqE,\tau)/\cqN.
\end{equation}
Here $\tau$ denotes precisely the source's topology. No equality of
this closed space with a principal algebraic ideal is being asserted.
Multiplication $M f(s)=sf(s)$ is continuous in the topological ring;
it preserves the dense generating subspace, and hence $\cqN$.
Its induced operator $S$ is the operator of Proposition 3.6.
We use the source's spectral assertion as an explicitly imported theorem:
\begin{equation}\label{eq:cq-source-spectrum}
 \cqSpec(S)=\{\lambda\in\cqC:
             \zeta(\tfrac12+i\lambda)=0
             \text{ at a nontrivial zero}\}.
\end{equation}
Spectrum here means failure of a continuous two-sided inverse.
No proof of the entire source spectral theorem is claimed as a new
result of this manuscript.

There is a material source limitation. Page 19 calls $\cqE$ the
Hadamard topological ring of entire functions of order at most one;
it does not state seminorms or a convergence criterion. The cited
Connes--Marcolli Proposition 2.24, printed pages 379--381,
proves the polynomial-image assertion and explicitly discusses
different function-space realizations \cite{connesMarcolli2008}. It does not supply
that missing specification. We therefore retain $\tau$ in every
source quotient. Below, all assertions depending on another topology
name and define that topology. The comparison with the source is an
exact quotient diagram, rather than an identification by terminology.

\subsection{A globally defined heat map on the original underlying functions}
\label{sec:cq-heatspace}
For an entire function $f$, let
$M_f(r)=\max_{|s|\leq r}|f(s)|$ and define its order by
\[
 \limsup_{r\to\infty}
       \frac{\log\log\max(e,M_f(r))}{\log r}.
\]
The zero function is included. For $1<p<2$, $a>0$ set
\[
 \|f\|_{p,a}=\sup_{s\in\cqC}|f(s)|e^{-a|s|^p}.
\]
The underlying set $\cqE$ is exactly the set where all these seminorms
are finite. Indeed order at most one gives, for $1<q<p$,
$M_f(r)\leq\exp(r^q)$ for sufficiently large $r$; the expression
$r^q-ar^p$ is bounded above. Conversely finiteness gives order at
most $p$ for every $p>1$.
Write $\gamma$ for this explicitly defined growth topology.
It is enough to take $p=1+1/j$, $j\geq2$, and $a=1/k$, $k\geq1$:
given $p>1,a>0$, select one of these $p_j<p$ and use the boundedness
of $|s|^{p_j}/k-a|s|^p$. This proves equivalence of the countable
family with all displayed seminorms. A sequence Cauchy in these
seminorms converges uniformly on compact sets to an entire function.
Taking pointwise limits in each weighted Cauchy inequality gives
convergence in that seminorm and finiteness of its norm. Thus
$(\cqE,\gamma)$ is a complete metrizable locally convex space.

\begin{cqlemma}[Derivative bounds and the heat group]\label{thm:cq-heat}
Put $D=\cqD$. Differentiation and multiplication by $s$ map the
original underlying set $\cqE$ into itself. For $m\geq1$,
\begin{equation}\label{eq:cq-cauchy}
 \|D^mf\|_{p,a}\leq
 m!\left(\frac{epa}{m}\right)^{m/p}
 \|f\|_{p,a/2^{p-1}}.
\end{equation}
For every $t\in\cqC$, the series
\begin{equation}\label{eq:cq-U}
 U_tf=\sum_{k=0}^\infty\frac{(-t/4)^k}{k!}D^{2k}f
\end{equation}
converges in $\gamma$, locally uniformly in $t$; $U_t$ is a
continuous automorphism of $(\cqE,\gamma)$ and
$U_tU_u=U_{t+u}$, $U_t^{-1}=U_{-t}$.
\end{cqlemma}
\begin{proof}
Cauchy's circle formula on a circle of radius $R$ about $s$ gives
\[
 |D^mf(s)|\leq m!R^{-m}\max_{|z-s|=R}|f(z)|.
\]
For $b=a/2^{p-1}$ the inequality
$|s+z|^p\leq2^{p-1}(|s|^p+|z|^p)$ bounds the last expression,
after multiplication by $e^{-a|s|^p}$, by
$m!R^{-m}e^{aR^p}\|f\|_{p,b}$. Choosing
$R=(m/(ap))^{1/p}$ proves \eqref{eq:cq-cauchy}.
For multiplication use
\[
 \|Mf\|_{p,a}\leq\sup_{r\geq0}re^{-ar^p/2}\|f\|_{p,a/2}.
\]
The $k$th heat summand is at most
\[
 \frac{(|t|/4)^k}{k!}(2k)!
       \left(\frac{epa}{2k}\right)^{2k/p}\|f\|_{p,b}.
\]
Since $(2k)!/k!\leq(2k)^k$, its scalar factor is at most
$[C_{p,a}|t|k^{1-2/p}]^k$, whose $k$th root tends to zero.
This proves summability, continuity, and locally uniform convergence
of every time-differentiated series by the same estimate with the
additional polynomial factors in $k$. Commutation with $D$ follows
from continuity of $D$. In the double composition series group
terms with $n=k+\ell$. The sum of the absolute scalar coefficients
at $D^{2n}f$ is $( (|t|+|u|)/4)^n/n!$; the preceding estimate
justifies regrouping. The binomial theorem gives $U_{t+u}$.
The value at zero is the identity, proving the inverse.
\end{proof}

\begin{cqlemma}[Polynomial approximation]\label{thm:cq-density}
The Taylor polynomials of any $f\in\cqE$ converge to $f$ in $\gamma$.
Multiplication of two elements is continuous in $\gamma$.
\end{cqlemma}
\begin{proof}
If $f=\sum c_ns^n$, Cauchy's coefficient bound with parameter
$b<a$ gives, for $n\geq1$,
\[
 |c_n|\leq\|f\|_{p,b}(epb/n)^{n/p}.
\]
Also $\|s^n\|_{p,a}=(n/(epa))^{n/p}$. Therefore
$\|c_ns^n\|_{p,a}\leq\|f\|_{p,b}(b/a)^{n/p}$.
The tail is bounded by a convergent geometric series.
Finally $\|fg\|_{p,a}\leq\|f\|_{p,a/2}\|g\|_{p,a/2}$.
\end{proof}

Retain the actual real heat family, rather than substitute a polynomial:
\begin{align}\label{eq:cq-realheat}
 \Phi(u)&=\sum_{n\geq1}
 (2\pi^2n^4e^{9u}-3\pi n^2e^{5u})e^{-\pi n^2e^{4u}},\notag\\
 Z_t(s)&=8H_t(2s)=8\int_0^\infty e^{tu^2}\Phi(u)\cos(2su)\,du,
 &\partial_tZ_t=-\tfrac14D^2Z_t,
 &Z_0=\Xi.
\end{align}
These are the Rodgers--Tao conventions, Section 1 \cite{rodgersTao2020}.
For clarity the analytic exchanges in this transport can be proved
directly. Since $n^2e^{4u}\geq(n^2+e^{4u})/2$, the absolutely
convergent sums $\sum n^je^{-\pi n^2/2}$ give constants $C,c>0$
such that $|\Phi(u)|\leq Ce^{9u}e^{-ce^{4u}}$ for $u\geq0$.
For $|t|\leq T$ and $|s|\leq R$, every derivative of the integrand
is bounded by a polynomial in $u$ times
$C\exp(Tu^2+(2R+9)u-ce^{4u})$, which is integrable.
The exponential series in $t$ can therefore be integrated absolutely.
Because $D^{2k}\cos(2su)=(-1)^k(2u)^{2k}\cos(2su)$, it gives
exactly $Z_t=U_t\Xi$ with the coefficient $(-t/4)^k/k!$.
In particular $Z_t$ lies in the original order-at-most-one set for
every complex $t$, by Lemma~\ref{thm:cq-heat}; it is never identically
zero because $U_t$ is invertible and $\Xi(0)\ne0$.

\subsection{Transport of the actual closed arithmetic image}
\label{sec:cq-moving}
Define $\tau_t$ on the same underlying vector space $\cqE$ by
declaring $U_t:(\cqE,\tau)\longrightarrow(\cqE,\tau_t)$ a
homeomorphism. This definition specifies every open set:
$V$ is $\tau_t$-open exactly when $U_{-t}V$ is $\tau$-open.
It uses the globally defined bijection just proved and does not require
unproved continuity in $\tau$. Set
\begin{equation}\label{eq:cq-Qt}
 N_t=U_t\cqN
 =\overline{\operatorname{span}\{U_t(R_\ell^\pm\Xi)\}}^{\,\tau_t},
 \qquad Q_t=(\cqE,\tau_t)/N_t.
\end{equation}
The closure equality follows because homeomorphisms carry closure to
closure. In particular $N_t$ is closed. The quotient map
\begin{equation}\label{eq:cq-Phi}
 \mathcal U_t:\cqQ\longrightarrow Q_t,\qquad [f]\longmapsto[U_tf]
\end{equation}
is well defined, bijective, and continuous, with continuous inverse
$[g]\mapsto[U_{-t}g]$. Indeed changing $f$ by $n\in\cqN$
changes $U_tf$ by $U_tn\in N_t$, and the inverse gives necessity;
the quotient universal property proves both continuity statements.

\begin{cqtheorem}[Exact arithmetic heat intertwiner]\label{thm:cq-moving}
On the entire-function set,
\begin{equation}\label{eq:cq-B}
 B_t=U_tMU_{-t}=M-\tfrac t2D.
\end{equation}
This operator preserves $N_t$, is continuous in $\tau_t$, and induces
$S_t$ on $Q_t$. It satisfies
$S_t\mathcal U_t=\mathcal U_tS$ and
\begin{equation}\label{eq:cq-specmove}
 \cqSpec(S_t)=\cqSpec(S)
 =\{\lambda:\zeta(\tfrac12+i\lambda)=0\text{ nontrivially}\}.
\end{equation}
The time-dependent trace functions are exactly
$U_t(R_\ell^\pm\Xi)=R_\ell^\pm(B_t)Z_t$.
\end{cqtheorem}
\begin{proof}
The product rule gives $D^{2k}Mf=MD^{2k}f+2kD^{2k-1}f$.
Insert this identity in the convergent heat series to obtain
$U_tM=(M-(t/2)D)U_t$, proving \eqref{eq:cq-B}.
If $n\in\cqN$, then $B_tU_tn=U_tMn\in N_t$.
Continuity is conjugation of the continuous $M$ by the defining
homeomorphism. The induced identity follows on representatives.
If $T$ is a continuous inverse of $S-\lambda$, then
$\mathcal U_tT\mathcal U_t^{-1}$ is a continuous inverse of
$S_t-\lambda$. Conjugating in the opposite direction proves the
converse, and hence the spectrum equality using
\eqref{eq:cq-source-spectrum}. Apply the polynomial identity
$U_tp(M)=p(B_t)U_t$ to $\Xi$ for the last assertion.
\end{proof}

The heat map also transports the full multiplication, with every
cross term retained. For $f,g\in\cqE$ define
\begin{equation}\label{eq:cq-star}
 f\star_t g=U_t\bigl((U_{-t}f)(U_{-t}g)\bigr)
    =\sum_{n=0}^\infty\frac{(-t/2)^n}{n!}(D^nf)(D^ng).
\end{equation}
We prove the series identity, including convergence. For
$F_t=U_tf$, $G_t=U_tg$, the series
\[
 J_t=\sum_{n=0}^\infty\frac{(-t/2)^n}{n!}
                    (D^nF_t)(D^nG_t)
\]
converges in $\gamma$, uniformly on compact time sets: applying
\eqref{eq:cq-cauchy} with output norm $(p,a/2)$ to both factors
bounds its $n$th term by
\[
 K\left[\frac T2\left(\frac{epa}{2}\right)^{2/p}
                   n^{1-2/p}\right]^n\quad(|t|\leq T),
\]
where $K$ bounds the product of the two input norms
$(p,a/2^p)$, and is finite by the heat estimate. The root test
applies. Cauchy's formula in time permits termwise time derivatives;
continuity of $D$ permits spatial derivatives. Differentiating the
coefficient contributes
\[
 -\tfrac12\sum_{n\ge0}(-t/2)^n(D^{n+1}F_t)(D^{n+1}G_t)/n!.
\]
Differentiating the factors contributes the two pure second-derivative
terms with coefficient $-1/4$. The complete product rule therefore
gives $\partial_tJ_t=-D^2J_t/4$ and $J_0=fg$.
Its entire time Taylor coefficients are uniquely
$(-1/4)^kD^{2k}(fg)/k!$, so $J_t=U_t(fg)$.
Replacing the inputs by $U_{-t}f,U_{-t}g$ proves
\eqref{eq:cq-star}. Conjugation of ordinary multiplication proves
associativity, commutativity, unit $1$, and continuity for $\tau_t$,
even when continuity of the series in the source topology is unknown.
The identity $s\star_t f=sf-tf'/2$ recovers exactly $B_t$.
No ideal property of $\cqN$ beyond its stated polynomial invariance
is required for this identity or for the quotient in
Theorem~\ref{thm:cq-moving}.

The transported actual trace image has no common point-evaluation
zero at any nonzero time. Suppose $\lambda$ were such a common
zero. It would annihilate each $U_t(s^n\Xi)$. For every $a\in\cqC$,
the Taylor series of $e^{as}$ converges in $\gamma$ by
Lemma~\ref{thm:cq-density}. Multiplication, $U_t$, and evaluation
are continuous in $\gamma$, so
\[
 0=\sum_{n=0}^\infty\frac{a^n}{n!}U_t(s^n\Xi)(\lambda)
   =e^{a\lambda-ta^2/4}Z_t(\lambda-ta/2).
\]
The last equality follows from the just-proved product formula
and the entire Taylor formula for $Z_t$. At $t\ne0$ varying $a$
forces $Z_t$ to vanish everywhere, a contradiction. This argument
uses only the polynomial generators contained in $N_t$; it assumes
no source-topology convergence of the exponential series. Thus the
source arithmetic spectrum in \eqref{eq:cq-specmove} is carried by
the explicitly transported operator and relation, even though the
transported relation has no common ordinary evaluation zero.

The first two nonconstant polynomial generators make the extra terms
visible without changing any polynomial coefficient:
\begin{align}\label{eq:cq-generators}
 U_t(s\Xi)&=sZ_t-\tfrac t2Z_t',\notag\\
 U_t(s^2\Xi)&=s^2Z_t-t sZ_t'-\tfrac t2Z_t+\tfrac{t^2}4Z_t''.
\end{align}
For $L=\partial_t+\tfrac14D^2$, the complete product rule is
\begin{equation}\label{eq:cq-unit}
 L(AP)=A\bigl(P_t+\tfrac14P_{ss}\bigr)
       +\bigl(A_t+\tfrac14A_{ss}\bigr)P
       +\tfrac12A_sP_s.
\end{equation}
Thus $L(R(s)Z_t)=\tfrac14R''Z_t+\tfrac12R'Z_t'$.
The family $R(B_t)Z_t$ solves the homogeneous heat equation, whereas
the displayed expression is the full defect of replacing it by
$R(s)Z_t$. For a local factorization $Z_t=AP$ with $A\ne0$,
the same formula retains the analytic unit $A$ and its derivatives.
At a simple zero $\lambda(t)$, differentiating
$Z_t(\lambda(t))=0$ gives
\begin{equation}\label{eq:cq-zero}
 \lambda'(t)=\frac14\frac{Z_{ss}}{Z_s}(t,\lambda(t))
 =\frac14\frac{P_{ss}}{P_s}(t,\lambda(t))
    +\frac12\frac{A_s}{A}(t,\lambda(t)).
\end{equation}
This statement is on the open locus $Z_s\ne0$, on which the analytic
implicit-function theorem supplies the local curve; it claims no
chosen zero is simple. At a multiple zero the division is undefined,
and the full equation \eqref{eq:cq-unit} remains the equation to use.

\subsection{Derivative domains, multivalued operators, and jet extensions}
\label{sec:cq-jets}
Every derivative has the full underlying domain $\cqE$ by
Lemma~\ref{thm:cq-heat}. Continuity in the unnamed source topology is
not assumed. The following gives exact source-quotient operators and
their domains without imposing it. For $k\geq1$ put
\[
 K_k=\{n\in\cqN:D^kn\in\cqN\},\qquad
 W_k=(D^k\cqN+\cqN)/\cqN\subseteq\cqQ.
\]
The derivative correspondence is the actual linear relation
\begin{equation}\label{eq:cq-relation}
 \mathfrak D_k=\{([f],[D^kf]):f\in\cqE\}
                    \subseteq\cqQ\times\cqQ.
\end{equation}
For fixed $[f]$, its output fibre is precisely $[D^kf]+W_k$:
all representatives are $f+n$ with $n\in\cqN$, and linearity gives
all and only the stated outputs. The kernel of its presenting map
$f\mapsto([f],[D^kf])$ is exactly $K_k$. The relation has domain
all of $\cqQ$; its multivalued part at zero is exactly $W_k$.
Thus it is the graph of an operator exactly when $D^k\cqN\subseteq
\cqN$. This is a calculated equivalence, with an explicit obstruction
space, rather than an assumption of invariance.

There is no nonempty domain of quotient classes on which the same
representative formula is unambiguous if $W_k\ne0$ and every
representative of a class is admitted: addition of the same offending
$n$ gives two different outputs for every class. A chosen linear
lift is a different map. More precisely any vector subspace
$A\subseteq\cqE$ supplies a single-valued map on its quotient image
$q(A)$ exactly when $D^k(A\cap\cqN)\subseteq\cqN$; the identical
difference-of-representatives argument proves both directions.
The first heat generator is $-D^2/4$, so its ambiguity space is
$-W_2/4$, with the coefficient and sign unchanged.

Even when $W_k$ is not closed, the smallest source-closed enlargement
of the target relation subspace is explicit:
\begin{equation}\label{eq:cq-target}
 [f]\longmapsto[D^kf]:
 \cqQ\longrightarrow
 \cqE/\overline{\cqN+D^k\cqN}^{\,\tau}.
\end{equation}
It is well defined algebraically. Equip its source with the quotient
of the graph topology $f\mapsto(f,D^kf)$ in
$(\cqE,\tau)^2$ to obtain a continuous map. This specifies a
topology change explicitly, without claiming it is the original one.
The lost part of the original target $\cqQ$ is exactly
$\overline{\cqN+D^k\cqN}^{\,\tau}/\cqN$.

Here is an extension which retains correlated derivatives rather than
quotienting them away. For $j\geq1$ define the bijection
\[
 T_j(f_0,\ldots,f_j)
   =(f_0,f_1-Df_0,\ldots,f_j-D^jf_0).
\]
Its inverse is $(f_0,g_1,\ldots,g_j)\mapsto
(f_0,g_1+Df_0,\ldots,g_j+D^jf_0)$.
Give $V_j=\cqE^{j+1}$ the topology transported by $T_j$ from
$(\cqE,\tau)\times(\cqE,\gamma)^j$. The residual factors
are explicitly given the proved Hausdorff growth topology; the
source factor keeps its original topology. For $J_jf=(f,Df,\ldots,D^jf)$,
$T_j(J_j\cqN)=\cqN\oplus0^j$, so $J_j\cqN$ is closed.
Define
\begin{equation}\label{eq:cq-jetextension}
 X_j=V_j/J_j\cqN.
\end{equation}
Then the map
\[
 [f_0,\ldots,f_j]\longmapsto
 ([f_0],f_1-Df_0,\ldots,f_j-D^jf_0)
\]
is a homeomorphism $X_j\simeq\cqQ\oplus(\cqE,\gamma)^j$; independence,
the inverse, and continuity follow directly from $T_j$ and the
quotient definition. In particular there is the split exact sequence
\begin{equation}\label{eq:cq-jetseq}
 0\longrightarrow\cqE^j\longrightarrow X_j
 \xrightarrow{[f_0,\ldots,f_j]\mapsto[f_0]}\cqQ
 \longrightarrow0,
 \qquad [f]\longmapsto[J_jf].
\end{equation}
Its relation on a source generator retains every product term:
\begin{equation}\label{eq:cq-rawjet}
 D^k(R\Xi)=\sum_{h=0}^k\binom kh R^{(h)}\Xi^{(k-h)}
                    \quad(0\leq k\leq j).
\end{equation}
The binomial formula follows by induction from the ordinary product
rule, including the two edge terms at each step.

The original spectral coordinate acts on $V_j$ by
\[
 (\mathcal M_j f)_k=s f_k+k f_{k-1}\quad(0\leq k\leq j),
 \qquad f_{-1}=0.
\]
The identity $D^k(sf)=sD^kf+kD^{k-1}f$ proves
$\mathcal M_jJ_jf=J_j(sf)$; it therefore preserves the exact
relation subspace. In the $T_j$ coordinates it acts by $S$ on
$\cqQ$ and by $g_k\mapsto sg_k+kg_{k-1}$ on the residual terms
($g_0=0$). This also proves continuity in the stated topology.
Notice that \eqref{eq:cq-jetseq} adds a free residual module:
the spectrum of the entire extension is not claimed to equal the
source spectrum. The arithmetic quotient and its split inclusion
are explicitly preserved.

\subsection{A fully computed analytic realization and its exact comparison}
\label{sec:cq-analytic}
We now compute the closures and the failure of descent in specified
topologies. This does not identify them with $\tau$.
Let $\kappa$ be the topology of uniform convergence on compact sets
on the same set $\cqE$, and define
\[
 N_\gamma=\overline{\cqC[s]\Xi}^{\,\gamma},\qquad
 N_\kappa=\overline{\cqC[s]\Xi}^{\,\kappa}.
\]
Evaluations of all derivatives are continuous in both topologies by
Cauchy's formula. The growth topology is finer than $\kappa$, so
$N_\gamma\subseteq N_\kappa$.
If $\lambda$ is a zero of $\Xi$ with its actual multiplicity $m_\lambda$,
every element of either closed subspace has vanishing derivatives
of orders $0,\ldots,m_\lambda-1$ at $\lambda$.

\begin{cqtheorem}[Compact closure, multiplicities, and coordinate spectrum]
\label{thm:cq-compact}
The compact closure is exactly
\begin{equation}\label{eq:cq-compactideal}
 \begin{aligned}
 N_\kappa
 &=\{f\in\cqE:f/\Xi\text{ extends to an entire function}\}\\
 &=\{f\in\cqE:D^hf(\lambda)=0\ (h<m_\lambda),\ \Xi(\lambda)=0\}.
 \end{aligned}
\end{equation}
On $Q_\kappa=(\cqE,\kappa)/N_\kappa$, multiplication by $s$ has
spectrum exactly the zeros of $\Xi$ in the original $s$ coordinate.
For a zero $\lambda$ of multiplicity $m$, there is a continuous
surjection
\[
 Q_\kappa\longrightarrow\cqC[\epsilon]/(\epsilon^m),\qquad
 [f]\longmapsto\sum_{h=0}^{m-1}\frac{D^hf(\lambda)}{h!}\epsilon^h,
\]
intertwining multiplication by $s$ with $\lambda+\epsilon$.
\end{cqtheorem}
\begin{proof}
The vanishing conditions are closed and hold on every $p\Xi$,
so the closure is contained in the last displayed set. Local Taylor
division at each isolated zero identifies that set with entire
divisibility by $\Xi$. For such an $f$, put $g=f/\Xi\in\cqO(\cqC)$.
Its Taylor polynomials $p_n$ converge uniformly on every compact;
$p_n\Xi\to f$ in $\kappa$. Every $p_n\Xi$ belongs to the original
order class even though no growth assertion about $g$ is needed.
This proves \eqref{eq:cq-compactideal} with the closure retained.
The local-jet map is independent of representatives; the polynomials
$1,(s-\lambda),\ldots,(s-\lambda)^{m-1}$ prove surjectivity.
Leibniz's rule gives the stated action, and the finite jet evaluations
are continuous.

If $\Xi(\lambda)\ne0$, define
\[
 T_\lambda f(s)=
 \frac{f(s)-f(\lambda)\Xi(s)/\Xi(\lambda)}{s-\lambda},
\]
filling in the removable value. Division by a linear factor preserves
the order bound: outside a disk about $\lambda$ the denominator is
bounded below, and inside a disk the filled entire function is bounded.
It is continuous in $\kappa$ by Cauchy's formula on a larger disk
applied to the numerator. For $n\in N_\kappa$ the quotient
$T_\lambda n/\Xi=(n/\Xi-(n/\Xi)(\lambda))/(s-\lambda)$ is entire;
thus it descends. Direct multiplication proves
$(S_\kappa-\lambda)[T_\lambda f]=[f]$.
Applying $T_\lambda$ to $(s-\lambda)f$ gives $f$ exactly, proving
the other inverse identity. If $\Xi(\lambda)=0$, evaluation at
$\lambda$ descends to a nonzero continuous functional which annihilates
the image of $S_\kappa-\lambda$. Hence that operator is not
surjective. This proves the spectrum claim.
\end{proof}

\begin{cqtheorem}[Direct differential and heat obstructions]
\label{thm:cq-obstruction}
Neither $N_\gamma$ nor $N_\kappa$ is invariant under $D$ or $D^2$.
For either frozen quotient, neither representative formula $[f']$
nor $[-f''/4]$ defines an operator on any nonempty class domain
admitting all representatives.
For every $t\ne0$, neither $U_tN_\gamma\subseteq N_\gamma$ nor
$U_tN_\kappa\subseteq N_\kappa$ holds. Ordinary multiplication
by $s$ preserves neither $U_tN_\gamma$ nor $U_tN_\kappa$.
\end{cqtheorem}
\begin{proof}
Choose any zero $\lambda$ of the actual $\Xi$ and let $m\geq1$
be its multiplicity. The existence of its zeros is part of the
classical arithmetic datum in the primary source; no realness or
simplicity is used. Write $\Xi=(s-\lambda)^mA$ with
$A(\lambda)\ne0$. Then $\Xi'=(s-\lambda)^{m-1}
(mA+(s-\lambda)A')$ does not vanish to order $m$.
Since $\Xi\in N_\gamma\subseteq N_\kappa$, both derivative
invariance assertions fail. If $m\geq2$, $\Xi''$ has leading
coefficient $m(m-1)A(\lambda)$ at order $m-2$, proving the
second-derivative failure. If $m=1$, use the actual polynomial
multiple $n=(s-\lambda)\Xi$ instead: $n''(\lambda)=2\Xi'(\lambda)
\ne0$. This retains the product-rule cross term and covers all
multiplicities. The nonempty-domain assertion follows from the
explicit fibre computation in \eqref{eq:cq-relation}.

For $a\in\cqC$, $e^{as}\Xi\in N_\gamma$ by polynomial
approximation of $e^{as}$ and continuous multiplication, and hence
belongs to $N_\kappa$. The conjugation $D(e^{as}f)=e^{as}(D+a)f$
and the absolutely convergent heat and translation series give
\begin{equation}\label{eq:cq-exponential}
 U_t(e^{as}\Xi)(s)
    =e^{as-ta^2/4}Z_t(s-ta/2).
\end{equation}
For completeness, the translation series for $f$ is its Taylor
series about $s$ and sums to $f(s-ta/2)$; its local uniform
convergence and the derivative estimates justify commutation of
$e^{-(ta/2)D}$ with $U_t$. If $U_tN_\bullet\subseteq N_\bullet$
for either choice, evaluate \eqref{eq:cq-exponential} at $\lambda$.
The exponential never vanishes, so $Z_t(\lambda-ta/2)=0$ for
every $a$. When $t\ne0$ these arguments fill $\cqC$; this contradicts
$Z_t\ne0$. Finally if $M$ preserved $U_tN_\bullet$, conjugation
would imply $(M+(t/2)D)N_\bullet\subseteq N_\bullet$.
Since $M$ already preserves $N_\bullet$, this forces
$DN_\bullet\subseteq N_\bullet$, contrary to the first part.
\end{proof}

The analytic moving-divisor quotient, which does carry the zeros of
the actual $Z_t$ under ordinary $s$, is
\begin{equation}\label{eq:cq-divisor}
 I_t=\overline{\cqC[s]Z_t}^{\,\kappa}
 =\{f\in\cqE:f/Z_t\text{ is entire}\},\qquad
 A_t=(\cqE,\kappa)/I_t.
\end{equation}
The proof of Theorem~\ref{thm:cq-compact} applies verbatim with
the nonzero entire order-at-most-one function $Z_t$: it uses only
local Taylor division and compact polynomial approximation.
Consequently $\cqSpec(M\text{ on }A_t)=\{\lambda:Z_t(\lambda)=0\}$,
with the full local multiplicity jets, and with the explicit inverse
$[f]\mapsto[(f-f(\lambda)Z_t/Z_t(\lambda))/(s-\lambda)]$
off that set. Equation \eqref{eq:cq-zero} is the evolution of this
original-coordinate spectrum on its simple-zero locus.
This spectrum can move; \eqref{eq:cq-specmove} is the spectrum
of the different, exactly specified transported operator $B_t$.

We now give the promised comparison morphisms with every kernel
retained. Put $J=\cqN\cap N_\kappa$ and
$C=(\cqE,\tau\vee\kappa)/J$, where $\tau\vee\kappa$ is the
topology generated by both topologies. Since each summand is closed
in its own topology, $J$ is closed in their join. The identity maps
give continuous surjections
\begin{equation}\label{eq:cq-span}
 \cqQ\ \longleftarrow\ C\ \longrightarrow\ Q_\kappa,
 \qquad\ker(C\to\cqQ)=\cqN/J,
 \quad\ker(C\to Q_\kappa)=N_\kappa/J.
\end{equation}
Each assertion follows by taking the same representative $f$ and
testing membership in the respective relation subspace; the quotient
universal property proves continuity. An algebraic common quotient is
$\cqE/(\cqN+N_\kappa)$ with the quotient topology from the join;
its Hausdorff version uses
$\overline{\cqN+N_\kappa}^{\,\tau\vee\kappa}$.
The map from $C$ to this further quotient is continuous. Continuity
of maps to it from each of the two separately topologized quotients
is not asserted. No closedness of the sum is presumed. This describes
the precise information lost in that further quotient.

At time $t$ replace the two subspaces by $N_t$ and $I_t$ and the
join by $\tau_t\vee\kappa$. The same construction gives continuous
surjections from their intersection quotient to the genuine transported
arithmetic quotient $Q_t$ and to $A_t$, with kernels
$N_t/(N_t\cap I_t)$ and $I_t/(N_t\cap I_t)$, respectively.
The two generator systems are linked by the explicit formula
$R(B_t)Z_t$ versus $R(s)Z_t$ and the defects
\eqref{eq:cq-generators}--\eqref{eq:cq-unit}.
These are exact morphisms and closure formulas for the actual source
space. Determining that either kernel vanishes for $\tau$ is not
established by the two inspected primary texts. Accordingly the
source-frozen statements $D\cqN\subseteq\cqN$ and
$D^2\cqN\subseteq\cqN$ have not been decided solely from that
unnamed topology. The unqualified noninvariance results above have
been proved in the two explicitly specified analytic realizations.

\subsection{The inverse-root cover, the branch locus, and arithmetic fibres}
\label{sec:cq-cover}
The retained polynomial inverse-fibre slice is
$P_{a,0,0}(r)=-2r^2-2a$. Its root equation is $a=-r^2$.
The arithmetic coordinate is $a=2s$, so
\begin{equation}\label{eq:cq-cover}
 s=-r^2/2.
\end{equation}
The finite free extension of the actual quotient is
\[
 \widehat Q=\cqC[r]\otimes_{\cqC[s]}\cqQ
           =\cqQ\oplus r\cqQ,
 \qquad
 \mathcal R=\begin{pmatrix}0&-2S\\1&0\end{pmatrix},
 \qquad -\tfrac12\mathcal R^2=\begin{pmatrix}S&0\\0&S\end{pmatrix}.
\]
The even--odd decomposition of a polynomial in $r$ proves freeness
with basis $(1,r)$ and uniqueness coefficient by coefficient.
Multiplying $q_0+rq_1$ by $r$ proves the displayed matrices;
the finite direct-sum topology makes every induced map continuous.
The involution is $(q_0,q_1)\mapsto(q_0,-q_1)$; the trace is
$\operatorname{Tr}(q_0+rq_1)=2q_0$, with kernel $r\cqQ$.
It is $\cqC[s]$-linear. It does not provide a $\cqC[r]$-module
quotient with this kernel unless $S=0$: multiplying $rq$ by $r$
gives $-2Sq$ in the even summand. This proves the exact limitation,
rather than suppressing the odd component. The same construction
at time $t$ uses $S_t$ and is conjugated by
$\operatorname{diag}(\mathcal U_t,\mathcal U_t)$.

The original $s$-spectrum is unchanged by the extension: an inverse
of $S-\lambda$ duplicates to a diagonal inverse, and conversely
projection of an inverse of the doubled matrix to its first component,
following inclusion of that component, gives an inverse of
$S-\lambda$. Both directions are continuous. Moreover
\begin{equation}\label{eq:cq-rspec}
 \cqSpec(\mathcal R)
 =\{\rho\in\cqC:-\rho^2/2\in\cqSpec(S)\}.
\end{equation}
To prove this without an unsupported spectral-mapping citation,
solve $(\mathcal R-\rho)(u,v)=(a,b)$. The second row gives
$u=b+\rho v$ and the first becomes
$-2(S+\rho^2/2)v=a+\rho b$. A continuous inverse of
$S+\rho^2/2$ therefore gives a continuous inverse matrix.
Conversely the right-hand side $(a,0)$ and the second projection
of an inverse matrix supply an inverse of $-2(S+\rho^2/2)$,
with both identities checked by the same two rows.

Before any quotient by $\Xi$, let $\mathcal O_s$ be the sheaf of
holomorphic functions in $s$ and put
$B=\mathcal O_s[r]/(r^2+2s)$, identified with
$\mathcal O_s\oplus r\mathcal O_s$. On the punctured base $s\ne0$,
implicit differentiation gives
$dr/ds=-1/r=r/(2s)$. Thus the actual lifted derivative is
\begin{equation}\label{eq:cq-connection}
 \nabla_s\binom{f_0}{f_1}
 =\binom{f_0'}{f_1'+f_1/(2s)},\quad
 R=\begin{pmatrix}0&-2s\\1&0\end{pmatrix},\quad
 [\nabla_s,R]=-R^{-1}.
\end{equation}
Indeed with $A=\operatorname{diag}(0,1/(2s))$,
$R'+AR-RA=\left(\begin{smallmatrix}0&-1\\1/(2s)&0\end{smallmatrix}\right)
=-R^{-1}$. A second application gives
\begin{equation}\label{eq:cq-connection2}
 \nabla_s^2=\operatorname{diag}
 \left(D^2,D^2+s^{-1}D-\frac1{4s^2}\right).
\end{equation}
Equivalently, on a function of $r$ the derivative is
$-r^{-1}\partial_r$ and the heat generator is
$-\tfrac14r^{-2}\partial_r^2+\tfrac14r^{-3}\partial_r$.
Squaring $-r^{-1}\partial_r$ by the product rule proves this
with both signs preserved.

The branch extension is exact. The meromorphic connection sends
$B$ into $s^{-1}B$ and the connection one-form has a simple pole
with residue $\operatorname{diag}(0,1/2)$. Its maximal domain for
holomorphic output at $s=0$ is
\begin{equation}\label{eq:cq-branchdomain}
 \operatorname{Dom}_{\rm hol}\nabla_s
       =\mathcal O_s\oplus r s\mathcal O_s.
\end{equation}
In fact the odd coefficient is $f_1'+f_1/(2s)$; its only possible
pole has coefficient $f_1(0)/2$, proving necessity and sufficiency.
For the square, write $f_1=\sum_{n\geq0}a_ns^n$. Its image has
coefficient $(n^2-1/4)a_ns^{n-2}$. The negative powers vanish
exactly when $a_0=a_1=0$, so
\[
 \operatorname{Dom}_{\rm hol}\nabla_s^2
       =\mathcal O_s\oplus r s^2\mathcal O_s.
\]
For $k$ iterates the odd coefficient is
$\prod_{h=0}^{k-1}(n+1/2-h)a_ns^{n-k}$, whose prefactor is
never zero for integral $n$. Thus its domain is
$\mathcal O_s\oplus r s^k\mathcal O_s$.
No derivation $B\to B$ extending $\partial_s$ exists at the branch:
derivation of $r^2+2s=0$ would give $2r\delta(r)+2=0$,
forcing $r$ to be a unit although it vanishes at the branch point.
The logarithmic derivative $s\nabla_s$
does extend to all of $B$, with matrix
$sD+\operatorname{diag}(0,1/2)$.

These are sheaf and representative calculations; they do not assert
that $\nabla_s$ descends through the frozen arithmetic relation.
For a holomorphic local multiplier $R_0(s)$ the relation vector
$(R_0\Xi,0)$ acquires $(R_0'\Xi+R_0\Xi',0)$, and the odd
relation $(0,R_0\Xi)$ acquires
$(0,R_0'\Xi+R_0\Xi'+R_0\Xi/(2s))$.
The second iterate includes the full $2R_0'\Xi'$ term, as well as
the odd $s^{-1}D-1/(4s^2)$ terms of
\eqref{eq:cq-connection2}. Jet extension
\eqref{eq:cq-jetextension}, or the moving closed quotient
\eqref{eq:cq-Qt}, retains these relations exactly.

Finally, the branch point is not an arithmetic spectral point at
$t=0$. The source gives $\Xi(0)\ne0$, so its Proposition 3.6
makes $S$ continuously invertible; then $\mathcal R$ is invertible
by \eqref{eq:cq-rspec}. In the analytic divisor realization the
stronger local statement follows from $\Xi$ being a holomorphic
unit on a disk about zero: the quotient stalk is zero there.
For real $t$, $Z_t(0)>0$ directly from the retained kernel:
\[
 2\pi^2n^4e^{9u}-3\pi n^2e^{5u}
 =\pi n^2e^{5u}(2\pi n^2e^{4u}-3)>0\quad(u\geq0),
\]
and the positive integral converges by the bound already proved.
Thus the actual real heat divisor never meets the cover branch
point $s=0$. On a simple zero curve its retained lift satisfies
$r^2=-2\lambda$, $r'=-\lambda'/r$; together with
\eqref{eq:cq-zero} these are the original-coordinate and cover
equations, including the unit contribution. A nonreal $r$ alone
does not specify whether $\lambda=-r^2/2$ is nonreal.

\subsection{Exact ingress from compact Weil tests}
\label{sec:cq-tests}
The accompanying retained Weil-test calculation supplies the concrete
space $\mathcal E=C_c^\infty(\mathbb R;\cqC)$ with coordinate
operator $T=-i\partial_u$ and Fourier transform
\[
 (\mathscr Ff)(s)=\int_{\mathbb R}f(u)e^{-ius}\,du.
\]
Here we give the complete heat ingress and its inverse domains;
none of the following claims requires identifying the source
quotient topology with a test-space topology. Integration by parts
gives $\mathscr F(Tf)=s\mathscr Ff$ with no endpoint term.
If $\operatorname{supp}f\subset[-L,L]$, then
\[
 |\mathscr Ff(s)|\leq\|f\|_1e^{L|\operatorname{Im}s|},
 \qquad
 \|\mathscr Ff\|_{p,a}\leq
 \|f\|_1\sup_{r\geq0}e^{Lr-ar^p}<\infty.
\]
Thus the Fourier transform maps the actual test space into the
original entire-function set, continuously in $\gamma$ on every
fixed-support smooth-function space and hence in its usual
inductive-limit topology.

Define, for all complex $t$,
\begin{equation}\label{eq:cq-testheat}
 (G_tf)(u)=e^{tu^2/4}f(u),\qquad
 U_t\mathscr Ff=\mathscr F G_tf,\qquad
 T_t=G_tTG_{-t}=T+\frac{it}{2}u.
\end{equation}
The multiplication $G_t$ is a continuous automorphism on each
fixed-support smooth-function space, with inverse $G_{-t}$.
Indeed each derivative of the smooth multiplier is bounded on the
fixed compact support, and Leibniz's finite sum bounds every output
derivative seminorm by finitely many input ones. The same argument
for the inverse proves the assertion. The exponential series and
every $u$-derivative converge uniformly on a fixed compact set,
locally uniformly in $t$; consequently the family is entire in time
in that test topology. In the Fourier integral,
$D^{2k}\mathscr Ff=\mathscr F((-1)^ku^{2k}f)$.
The compact support permits absolute interchange with the
exponential series, proving the middle identity with the exact
factor $1/4$. Finally differentiation of $e^{-tu^2/4}f$ gives
the displayed $T_t$, and
\[
 \mathscr F(T_tf)=B_t\mathscr Ff,
 \qquad D\mathscr Ff=\mathscr F(-iu f).
\]
Both equations follow directly from the integral, including their
signs. The test-space version of the retained cover is therefore
\[
 \mathcal R_t^{\rm test}
   =\begin{pmatrix}0&-2T_t\\1&0\end{pmatrix},\qquad
 (\mathcal R_t^{\rm test})^2=-2\operatorname{diag}(T_t,T_t).
\]
Componentwise Fourier transformation intertwines this matrix with
the representative matrix using $B_t$, and then with the actual
quotient matrix using $S_t$.

To retain every possible arithmetic relation in this ingress, put
\begin{equation}\label{eq:cq-testkernel}
 K=\{f\in\mathcal E:\mathscr Ff\in\cqN\},\qquad
 K_t=\{f\in\mathcal E:\mathscr Ff\in N_t\}.
\end{equation}
Then $K_t=G_tK$: by \eqref{eq:cq-testheat},
$\mathscr Ff\in U_t\cqN$ exactly when
$\mathscr F(G_{-t}f)\in\cqN$. Thus there are injective linear
maps onto their explicitly specified images
\[
 \mathcal E/K\longrightarrow\cqQ,\quad[f]\longmapsto[\mathscr Ff],
 \qquad
 \mathcal E/K_t\longrightarrow Q_t,\quad[f]\longmapsto[\mathscr Ff],
\]
and the square formed by $G_t$ and $\mathcal U_t$ commutes.
The kernel statements follow exactly from
\eqref{eq:cq-testkernel}; no claim that $K=0$ in the unnamed
source topology is inserted. To make the ingress continuous without
that identification, give $\mathcal E$ the join $\eta$ of its
usual test topology with the inverse-image topology of $\tau$
under $\mathscr F$. Then $K$ is closed and the first quotient map
is continuous. Give the time-$t$ test space the transported
topology $(G_t)_*\eta$; the second map and the commutative square
are continuous by \eqref{eq:cq-testheat}. This specifies the exact
topologies rather than claiming an unproved topological embedding
of a quotient image equipped with its subspace topology.

There is a fully explicit inverse-domain calculation in the same
test coordinates. For $\lambda\in\cqC$, $k\geq1$, the image is
\begin{equation}\label{eq:cq-testmoments}
 (T_t-\lambda)^k\mathcal E
 =\left\{g\in\mathcal E:
 \int_{\mathbb R}v^j e^{-tv^2/4-i\lambda v}g(v)\,dv=0
                  \quad(0\leq j<k)\right\}.
\end{equation}
On that precise image the inverse is
\begin{equation}\label{eq:cq-testinverse}
 [(T_t-\lambda)^{-k}g](u)
 =\frac{i^ke^{tu^2/4+i\lambda u}}{(k-1)!}
 \int_{-\infty}^{u}(u-v)^{k-1}
            e^{-tv^2/4-i\lambda v}g(v)\,dv.
\end{equation}
To prove both directions, write
$f=e^{tu^2/4+i\lambda u}h$.
Direct differentiation gives
\[
 (T_t-\lambda)f=e^{tu^2/4+i\lambda u}(-i h').
\]
Iteration reduces the inverse equation to a $k$th primitive, giving
\eqref{eq:cq-testinverse} because $(-i)^ki^k=1$.
For $g$ supported in $[a,b]$, that formula vanishes below $a$.
Above $b$, expansion of $(u-v)^{k-1}$ shows it vanishes identically
exactly when all $k$ moments in \eqref{eq:cq-testmoments} vanish;
the binomial coefficients and the nonzero exponential do not change
that equivalence. Its smoothness follows by repeated differentiation
of the integral; smoothness of $g$ with compact support gives
matching derivatives at both endpoints. Necessity for any compact
solution also follows by integrating $h^{(k)}$ against
$1,v,\ldots,v^{k-1}$ and integrating by parts $k$ times.
Uniqueness follows since a compactly supported solution of
$h^{(k)}=0$ is a polynomial of degree at most $k-1$ which vanishes
on a ray, and hence is zero. This proves the exact image, inverse,
kernel zero, and preservation of the original support interval.

In particular the inverse of $\mathcal R_t^{\rm test}$ is
\[
 (g_0,g_1)\longmapsto(g_1,-\tfrac12T_t^{-1}g_0),
 \qquad
 \operatorname{Dom}(\mathcal R_t^{\rm test})^{-1}
 =\left\{(g_0,g_1):\int e^{-tv^2/4}g_0(v)\,dv=0\right\}.
\]
Substitution into its two rows proves both inverse identities.
Thus the branch-locus obstruction in the concrete test realization
is a retained Gaussian-weighted zero-moment condition. The exact
arithmetic quotient instead has invertible $S_t$ and root operator,
as proved above. Their ingress map and its kernels explain precisely
how these statements coexist; one operator's inverse domain is not
silently assigned to the other realization.

\subsection{Proof scope and reproducibility}
The new source-topology assertions proved here are the complete
transported closure, homeomorphic quotient, conjugated spectral
operator, source-domain differential correspondence, correlated jet
extension, finite free cover, compact-test heat ingress with its
weighted-moment inverse domains, and exact comparison diagrams. The
explicit compact and growth realizations have proved frozen
differential and heat noninvariance; the compact realization also
has a complete divisor and spectrum calculation with multiplicities.
The missing identification of the paper's unnamed topology is
recorded explicitly and is not replaced by a conditional arithmetic
theorem. Nothing in these results establishes an off-critical zero,
a negative Weil value, or an arithmetic endpoint by analogy with
the inverse-fibre collision. The exact transport through that
collision cover and the derivative information it requires have
been calculated above.

The companion checker uses rational polynomial arithmetic to verify
the heat conjugation on finite polynomial inputs, the retained
product terms, the matrix commutator, the iterated branch-domain
coefficients, the original-coordinate relation, and sample finite
jet multiplication. These are reproducible checks of the displayed
identities, not a substitute for their written proofs. Source
hashes and exact reading locators are in the companion provenance
receipt; no numerical zero search or Lean run is part of this work.
